% !TeX program = lualatex
% Compile twice: lualatex 83.tex
% Standalone research notebook; original section preserved.
% Research additions: 27 September 2026. No external TeX input or BibTeX file.
\documentclass[11pt,a4paper]{article}
\usepackage[margin=24mm,headheight=14pt]{geometry}
\usepackage{fontspec}
\setmainfont{Latin Modern Roman}
\setsansfont{Latin Modern Sans}
\setmonofont{Latin Modern Mono}
\usepackage{amsmath,amssymb,mathtools}
\usepackage{unicode-math}
\setmathfont{Latin Modern Math}
\usepackage{microtype}
\usepackage{enumitem,xurl,needspace}
\usepackage[dvipsnames]{xcolor}
\usepackage{hyperref}
\hypersetup{unicode=true,colorlinks=true,linkcolor=MidnightBlue,urlcolor=MidnightBlue,
 pdftitle={Research notebook - Problem 83},pdfauthor={Open Problems Math},bookmarksnumbered=true}
\usepackage{bookmark}
\usepackage{titlesec}
\titleformat{\section}{\Large\bfseries\raggedright}{\thesection}{0.7em}{}
\usepackage{fancyhdr}
\pagestyle{fancy}\fancyhf{}
\fancyhead[L]{\small\sffamily Open Problems Math | Research notebook}
\fancyhead[R]{\small\sffamily Problem 83 | 27 September 2026}
\fancyfoot[C]{\thepage}
\setlength{\parindent}{0pt}
\setlength{\parskip}{5pt plus 1pt}
\setlength{\emergencystretch}{3em}
\setcounter{tocdepth}{1}
\setcounter{secnumdepth}{1}
\setlist[itemize]{leftmargin=3em,itemsep=5pt,topsep=4pt,parsep=0pt}
\allowdisplaybreaks
\numberwithin{equation}{section}
\newcommand{\N}{\mathbb{N}}
\newcommand{\Z}{\mathbb{Z}}
\newcommand{\Q}{\mathbb{Q}}
\newcommand{\R}{\mathbb{R}}
\newcommand{\C}{\mathbb{C}}
\newcommand{\F}{\mathbb{F}}
\newcommand{\A}{\mathbb{A}}
\newcommand{\PP}{\mathbb P}
\newcommand{\E}{\mathbb{E}}
\newcommand{\eps}{\varepsilon}
\newcommand{\dd}{\,\mathrm d}
\newcommand{\sref}[2]{\hyperlink{source:#1:#2}{[S#2]}}
\newcommand{\eref}[2]{\hyperlink{extra:#1:#2}{[E#2]}}
\newif\ifshowcatalogue
\showcataloguetrue
\newcommand{\cataloguescope}[1]{\ifshowcatalogue
 \paragraph{Exact scope in the supplied catalogue.}
 \begin{quote}\small #1\end{quote}\fi}
\begin{document}
\setcounter{section}{82}
\section{\texorpdfstring{Bateman-Horn conjecture}{Bateman-Horn conjecture}}\label{83}
\subsection{Definitions and mathematical statement}
For fixed distinct nonconstant irreducible $f_1,\ldots,f_r\in\Z[x]$ of positive leading coefficient, let $d_i=\deg f_i$ and let $N_p$ count roots of $\prod_i f_i$ modulo $p$. Admissibility means $N_p<p$ for every prime. The Bateman--Horn prediction is
\[
 \#\{1\le n\le x:f_i(n)\text{ all prime}\}
 \sim\frac{C(f)}{\prod_i d_i}\frac{x}{(\log x)^r},\qquad
 C(f)=\prod_p\frac{1-N_p/p}{(1-1/p)^r}.
\]
Equivalently use $C(f)\int^x dt/\prod_i\log f_i(t)$, starting where all values exceed one. The polynomials are fixed before $x$ tends to infinity.
\par\smallskip\textit{Formulation and concept references:} \sref{83}{1}.
\cataloguescope{Prove or disprove the Bateman-Horn asymptotic formula for the number of positive integers n up to x for which any fixed admissible finite family of distinct irreducible integer polynomials with positive leading coefficients are simultaneously prime.}
\subsection{Short English statement}
Do simultaneous prime values of fixed admissible polynomials follow the exact density predicted by Bateman and Horn?
% BEGIN RESEARCH_ADDITION ATTEMPT_83
\subsection{Research attempt}

\paragraph{Current frontier.}
The original survey develops the singular series, admissibility and fixed
polynomial convention \eref{83}{1}. Coefficient-averaged prime-value
results do not remove exceptional fixed tuples. Likewise the 2026 work
on sets of integers satisfying Bateman--Horn statistics constructs sets
with prescribed statistics; it does not identify those sets with the
actual primes \sref{83}{3}. The unverified full-proof claim in
\sref{83}{2} is not used. The attempt below tests whether exact finite
congruence calculations can be pushed to the prime-factor cutoff needed
for an asymptotic, and calibrates the obstruction on the solved linear case.

\paragraph{Attack route.}
The local factors suggest a product of residue-class densities. Chinese
remaindering proves such a product at any fixed cutoff, but does not justify
letting the cutoff grow at the same rate as the input. A weighted-prime
formulation isolates the exact analytic conclusion needed. The route
combines these two descriptions and then checks the interchange of limits
against $f(n)=n$, where the true answer is known independently.

\paragraph{Attempt: a prime-only weighted equivalent.}
Choose $n_0$ so that every $f_i(t)>1$ and is increasing for $t\geq n_0$.
Write
\[
 \vartheta_*(m)=\begin{cases}\log m,&m\text{ is prime},\\0,&\text{otherwise},\end{cases}
 \qquad g(t)=\prod_{i=1}^r\log f_i(t),
\]
\[
 W(x)=\sum_{n_0\leq n\leq x}\prod_i\vartheta_*(f_i(n)),
 \qquad P(x)=\#\{n_0\leq n\leq x:f_i(n)\text{ all prime}\}.
\]
If
\begin{equation}\label{eq83:weighted}
 W(x)=C(f)x+o(x),
\end{equation}
then the target asymptotic follows. Abel summation gives, up to fixed
initial terms,
\begin{equation}\label{eq83:abel}
 P(x)=\frac{W(x)}{g(x)}+
              \int_{n_0}^x W(t)\frac{g'(t)}{g(t)^2}\,dt.
\end{equation}
Here
$g(t)\sim(\prod_i d_i)(\log t)^r$ and
$g'(t)/g(t)=O(1/(t\log t))$. The main term on the right is
$C(f)\int_{n_0}^xdt/g(t)$ plus a fixed constant, since
$(t/g(t))'=1/g(t)-t g'(t)/g(t)^2$.
If $W(t)=C(f)t+E(t)$ with $E(t)=o(t)$, its first error is
$o(x/g(x))$. For the integral error split at a fixed large $T_0$ on which
$|E(t)|\leq\varepsilon t$; the tail is at most
$\varepsilon\int_{T_0}^x t g'(t)/g(t)^2dt=O(\varepsilon x/g(x))$.
Sending $\varepsilon\to0$ proves the desired error. Thus the degree factor
is exactly $\prod_i d_i$, not an omitted normalization.

Conversely the stated asymptotic for $P(x)$ implies
\eqref{eq83:weighted} by summation against $g(n)$: its leading endpoint
term is $C(f)x+o(x)$ and the derivative term is $O(x/\log x)+o(x)$.
This is an equivalent analytic target, not an established estimate.
Using $\vartheta_*$ rather than $\Lambda$ avoids an unproved assertion
that proper prime powers among arbitrary high-degree polynomial values
are negligible. The trivial count of all prime powers up to $O(x^d)$
is not automatically $o(x)$ when $d$ is large.

\paragraph{Attempt: exact local models at a fixed cutoff.}
Let $P_z=\prod_{p\leq z}p$ and
$A_z=\prod_{p\leq z}(1-1/p)^{-r}$. Define
\[
 R_z(n)=A_z\,\mathbf1_{\gcd(P_z,\prod_i f_i(n))=1}.
\]
By the Chinese remainder theorem its exact mean over one full period is
\begin{equation}\label{eq83:local}
 \frac1{P_z}\sum_{n=1}^{P_z}R_z(n)
 =A_z\prod_{p\leq z}(1-N_p/p)=:C_z(f).
\end{equation}
Counting full periods and a remainder yields
\begin{equation}\label{eq83:error}
 \frac1X\sum_{n\leq X}R_z(n)
       =C_z(f)+O\left(\frac{A_z P_z}{X}\right).
\end{equation}
For fixed $z$ this proves the local density limit. The error deliberately
retains the modulus size; omitting it would be the invalid uniformity
step. In particular, admissibility makes the finite products positive,
but positivity of every finite product does not on its own supply
\eqref{eq83:weighted}.

\paragraph{Adversarial calibration: the wrong constant at the prime cutoff.}
Take the admissible single polynomial $f(n)=n$. Then $N_p=1$ and
$C_z(f)=1$ for every $z$. Nevertheless, with $z=\sqrt X$, all integers
$1\leq n\leq X$ coprime to $P_z$ are exactly $1$ and primes exceeding
$\sqrt X$: a composite integer in that range has a prime divisor at most
$\sqrt X$. Hence
\[
 \frac1X\sum_{n\leq X}R_{\sqrt X}(n)
 =\frac{1+\pi(X)-\pi(\sqrt X)}X
           \prod_{p\leq\sqrt X}(1-1/p)^{-1}.
\]
The prime number theorem and Mertens' product theorem give
\begin{equation}\label{eq83:mismatch}
 \frac1X\sum_{n\leq X}R_{\sqrt X}(n)
       \longrightarrow\frac{e^\gamma}{2}\ne1,
\end{equation}
where $\gamma$ is Euler's constant. These established inputs are discussed
in \eref{83}{1} and proved in the Mertens exposition \eref{83}{2}.
For the inequality note $0<\gamma<\log2$, for example from the usual
integral bounds on the harmonic sequence; a numerical approximation is
not needed to distinguish the constants.

Thus the local fixed-$z$ limit followed by $z\to\infty$ gives one, while
the diagonal cutoff needed to recognize primes gives a different value.
This explicitly falsifies the proposed unrestricted interchange of limits,
even before nonlinear or simultaneous values enter. It is not a
counterexample to Bateman--Horn: the ordinary prime count still has its
correct asymptotic. The normalization error belongs to the overextended
independent-local sieve model.

\paragraph{Attempt: the bridge that would actually be needed.}
For general fixed polynomials a viable comparison must account for large
prime factors and parity after the small-prime congruence restrictions.
One possible sufficient form is a decomposition
\[
 \prod_i\vartheta_*(f_i(n))=R_{z(X)}(n)+\mathcal E_X(n)
\]
with a calibrated correction whose averaged value tends to
$C(f)-\lim X^{-1}\sum_{n\leq X}R_{z(X)}(n)$, whenever that latter limit
exists. Merely requiring $\sum\mathcal E_X=o(X)$ at an arbitrary deep
cutoff is already false in the linear calibration. A different normalization
or a genuine bilinear/parity analysis is needed, not just better bookkeeping
of finitely many forbidden residues.

This observation is a constraint on an attack, not a new error estimate.
For a small cutoff where \eqref{eq83:error} is effective, the rough values
still include many composites. For a cutoff recognizing primality, the
simple CRT error is useless and the conditional distribution is not the
product model. The regime between them is the unsolved analytic content.

\paragraph{Stress test and computation.}
The polynomials are fixed throughout; varying coefficients with $X$ would
change the target. Repeated or proportional polynomial issues cannot be
used to manufacture a local obstruction after the original admissibility
and distinctness conditions are imposed. The companion script checks
finite CRT means and counts prime values for two sample admissible tuples,
without interpreting a finite ratio as an asymptotic proof. The all-prime
weight avoids silently discarding prime-power terms. A proof of the full
asymptotic would also imply the qualitative simultaneous-prime assertion
of Section~54 of the original catalogue; that unresolved consequence is
not an input used here.

\paragraph{Outcome.}
\textbf{Outcome: Exploratory approach with unresolved gap.}

The weighted equivalent and fixed-cutoff local identities are proved, and
the diagonal-cutoff constant mismatch gives a rigorous countertest to a
natural but false limit interchange. No new prime-value asymptotic or
admissible counterexample is obtained. The missing task is a calibrated
large-factor and parity estimate for every fixed tuple, strong enough for
\eqref{eq83:weighted}; the CRT computation alone cannot supply it.

% Keep the local bibliography together rather than leave a source fragment.
\needspace{170mm}
% END RESEARCH_ADDITION ATTEMPT_83
\subsection{Sources}
\begin{itemize}\raggedright
\item[\textnormal{[S1]}]\hypertarget{source:83:1}{} S. L. Aletheia-Zomlefer, L. Fukshansky, and S. R. Garcia, Expositiones Mathematicae 38 (2020), 430-479.
\url{https://www.sciencedirect.com/science/article/am/pii/S0723086918301178}.
{\small\textit{Catalogue formulation source.}}
\item[\textnormal{[S2]}]\hypertarget{source:83:2}{} From Golomb to Bateman-Horn.
\url{https://www.preprints.org/manuscript/202512.2666}.
{\small\textit{Catalogue research/status reference; not a proof endorsement.}}
\item[\textnormal{[S3]}]\hypertarget{source:83:3}{} Sets of integers satisfying Bateman-Horn statistics.
\url{https://arxiv.org/abs/2605.01155}.
{\small\textit{Catalogue research/status reference; not a proof endorsement.}}
% BEGIN RESEARCH_ADDITION REFERENCES_83
\item[\textnormal{[E1]}]\hypertarget{extra:83:1}{} S. L. Aletheia-Zomlefer,
Lenny Fukshansky and Stephan Ramon Garcia, \emph{The Bateman--Horn Conjecture:
Heuristics, History, and Applications}, Expositiones Mathematicae 38 (2020),
430--479; arXiv:1807.08899, formulation and Section 5.3.
\url{https://arxiv.org/html/1807.08899}.
{\small\textit{Inspected 27 September 2026; fixed-tuple and prime-counting
conventions retained.}}
\item[\textnormal{[E2]}]\hypertarget{extra:83:2}{} Terence Tao,
\emph{Mertens' theorems}, author's mathematical exposition, 11 December 2013,
including the Mertens product formula.
\url{https://terrytao.wordpress.com/2013/12/11/mertens-theorems/}.
{\small\textit{Used for the classical product normalization in the solved
linear calibration, not as support for a nonlinear prime conjecture.}}
% END RESEARCH_ADDITION REFERENCES_83
\end{itemize}

\end{document}
